\documentclass[12pt,a4paper]{article}
\usepackage{mathwizards-ingles}
\title{%
  \vspace{-1cm}
  {\Huge\bfseries\textcolor{mes2}{Mes 2 · Semana 7}}\\[0.3cm]
  {\Large Futuro, Verbos Modales y Phrasal Verbs}\\[0.2cm]
  {\large\color{grissuave} Nivel A2 · Lectura Extensiva y Nuevos Patrones}\\[0.5cm]
  \rule{\textwidth}{1.5pt}
}
\author{}
\date{}

\begin{document}
\maketitle
\thispagestyle{empty}

\begin{cajanivel}
\textbf{Objetivo:} El estudiante consume textos más largos en lectura extensiva (Graded Readers A1/A2). Se introducen las dos formas de \textbf{futuro} (\textit{will} y \textit{going to}), los \textbf{verbos modales básicos} y los primeros \textbf{phrasal verbs} de alta frecuencia. También se inician las \textbf{collocations} como unidades léxicas.
\end{cajanivel}

% ═══════════════════════════════════════════════════════════════════════════════
\section{Gramática Emergente: Futuro con \textit{will}}
% ═══════════════════════════════════════════════════════════════════════════════

\begin{cajagrammar}[Future Simple: \textit{will} + verbo base]
Se usa para \textbf{predicciones, promesas y decisiones espontáneas}.

\medskip
\begin{tabularx}{\textwidth}{@{} l X @{}}
\toprule
\textbf{Estructura} & \textbf{Ejemplo} \\
\midrule
\textbf{Afirmativo:} Sujeto + will + V & \textit{I \textbf{will help} you.} / \textit{I'\textbf{ll help} you.} \\
\textbf{Negativo:} Sujeto + won't + V & \textit{She \textbf{won't come} tomorrow.} \\
\textbf{Pregunta:} Will + sujeto + V? & \textit{\textbf{Will} you \textbf{be} at the party?} \\
\bottomrule
\end{tabularx}

\medskip
\begin{cajaejemplo}[¿Cuándo usar \textit{will}?]
\begin{itemize}[nosep, leftmargin=0.8cm]
  \item \textbf{Predicciones:} \textit{I think it \textbf{will rain} tomorrow.}
  \item \textbf{Promesas:} \textit{I \textbf{will call} you tonight.}
  \item \textbf{Decisiones espontáneas:} \textit{I'm hungry. I\textbf{'ll have} a sandwich.}
  \item \textbf{Ofertas:} \textit{I\textbf{'ll carry} that for you.}
\end{itemize}
\end{cajaejemplo}
\end{cajagrammar}

% ═══════════════════════════════════════════════════════════════════════════════
\section{Gramática Emergente: Futuro con \textit{going to}}
% ═══════════════════════════════════════════════════════════════════════════════

\begin{cajagrammar}[Future with \textit{going to}: am/is/are + going to + verbo base]
Se usa para \textbf{planes e intenciones} y \textbf{predicciones con evidencia}.

\medskip
\begin{tabularx}{\textwidth}{@{} l X @{}}
\toprule
\textbf{Estructura} & \textbf{Ejemplo} \\
\midrule
\textbf{Afirmativo} & \textit{I'\textbf{m going to study} English tonight.} \\
\textbf{Negativo} & \textit{We'\textbf{re not going to travel} this year.} \\
\textbf{Pregunta} & \textit{\textbf{Are} you \textbf{going to watch} the movie?} \\
\bottomrule
\end{tabularx}

\medskip
\begin{cajaejemplo}[¿Cuándo usar \textit{going to}?]
\begin{itemize}[nosep, leftmargin=0.8cm]
  \item \textbf{Planes:} \textit{We\textbf{'re going to visit} grandma this weekend.} (ya decidido)
  \item \textbf{Evidencia:} \textit{Look at those clouds! It\textbf{'s going to rain}.} (hay evidencia)
\end{itemize}
\end{cajaejemplo}
\end{cajagrammar}

\begin{cajagrammar}[Contraste Rápido: \textit{will} vs. \textit{going to}]
\begin{tabularx}{\textwidth}{@{} l X X @{}}
\toprule
& \textbf{\textit{will}} & \textbf{\textit{going to}} \\
\midrule
Decisión & Espontánea (ahora mismo) & Planeada (de antes) \\
Ejemplo & \textit{I'll have coffee.} (acabo de decidir) & \textit{I'm going to have coffee.} (ya lo tenía pensado) \\
Predicción & Sin evidencia (opinión) & Con evidencia (lo veo venir) \\
\bottomrule
\end{tabularx}
\end{cajagrammar}

% ═══════════════════════════════════════════════════════════════════════════════
\section{Gramática Emergente: Verbos Modales Básicos}
% ═══════════════════════════════════════════════════════════════════════════════

\begin{cajagrammar}[Modal Verbs: can, must, should]
Los modales van \textbf{antes del verbo base} y \textbf{no cambian} con la persona (no hay ``-s'' en tercera persona).

\medskip
\begin{tabularx}{\textwidth}{@{} l l X @{}}
\toprule
\textbf{Modal} & \textbf{Significado} & \textbf{Ejemplo} \\
\midrule
\eng{can} & habilidad / permiso & \textit{I \textbf{can} swim.} / \textit{\textbf{Can} I sit here?} \\
\eng{can't} & incapacidad / prohibición & \textit{He \textbf{can't} drive.} / \textit{You \textbf{can't} park here.} \\
\eng{must} & obligación fuerte & \textit{You \textbf{must} wear a seatbelt.} \\
\eng{mustn't} & prohibición & \textit{You \textbf{mustn't} touch that.} \\
\eng{should} & consejo & \textit{You \textbf{should} study more.} \\
\eng{shouldn't} & desaconsejado & \textit{You \textbf{shouldn't} eat so much sugar.} \\
\bottomrule
\end{tabularx}

\medskip
\begin{cajaejemplo}[Nota importante]
$\checkmark$ \textit{She \textbf{can} speak English.} \quad $\times$ \textit{She \textbf{cans} speak English.}\\
Los modales \textbf{nunca} llevan ``-s'', ``to'' ni cambian de forma.
\end{cajaejemplo}
\end{cajagrammar}

% ═══════════════════════════════════════════════════════════════════════════════
\section{Chunks: Collocations de Alta Frecuencia}
% ═══════════════════════════════════════════════════════════════════════════════

\begin{cajachunk}[Common Collocations — Combinaciones Naturales]
En inglés, ciertas palabras \textbf{siempre van juntas}. No se dice ``do a photo'', se dice ``\textit{take} a photo''. Estas combinaciones se aprenden como \textbf{chunks}, no palabra por palabra.

\medskip
\begin{multicols}{2}
\begin{tabularx}{\linewidth}{@{} l X @{}}
\toprule
\textbf{Con \textit{make}} & \\
\midrule
\eng{make a mistake} & \esp{cometer un error} \\
\eng{make a decision} & \esp{tomar una decisión} \\
\eng{make breakfast} & \esp{preparar el desayuno} \\
\eng{make money} & \esp{ganar dinero} \\
\eng{make a noise} & \esp{hacer ruido} \\
\eng{make friends} & \esp{hacer amigos} \\
\bottomrule
\end{tabularx}

\columnbreak

\begin{tabularx}{\linewidth}{@{} l X @{}}
\toprule
\textbf{Con \textit{take} / \textit{do}} & \\
\midrule
\eng{take a photo} & \esp{tomar una foto} \\
\eng{take a shower} & \esp{ducharse} \\
\eng{take a break} & \esp{tomar un descanso} \\
\eng{do homework} & \esp{hacer tarea} \\
\eng{do exercise} & \esp{hacer ejercicio} \\
\eng{do the dishes} & \esp{lavar los platos} \\
\bottomrule
\end{tabularx}
\end{multicols}

\medskip
\begin{tabularx}{\textwidth}{@{} l X @{}}
\toprule
\textbf{Con \textit{have}} & \\
\midrule
\eng{have a good time} \esp{pasarla bien} & \eng{have a meeting} \esp{tener una reunión} \\
\eng{have a conversation} \esp{tener una conversación} & \eng{have a problem} \esp{tener un problema} \\
\bottomrule
\end{tabularx}
\end{cajachunk}

% ═══════════════════════════════════════════════════════════════════════════════
\section{Chunks: Phrasal Verbs de Alta Frecuencia}
% ═══════════════════════════════════════════════════════════════════════════════

\begin{cajachunk}[Essential Phrasal Verbs — Los 15 Más Comunes]
Un \textit{phrasal verb} es un verbo + partícula que cambia el significado original del verbo. Se aprenden como \textbf{unidades completas}, nunca por partes.

\medskip
\begin{tabularx}{\textwidth}{@{} l l X @{}}
\toprule
\textbf{Phrasal Verb} & \textbf{Significado} & \textbf{Ejemplo} \\
\midrule
\eng{look for} & buscar & \textit{I'm \textbf{looking for} my keys.} \\
\eng{look after} & cuidar & \textit{She \textbf{looks after} her sister.} \\
\eng{give up} & rendirse & \textit{Don't \textbf{give up}!} \\
\eng{turn on} & encender & \textit{\textbf{Turn on} the TV.} \\
\eng{turn off} & apagar & \textit{Please \textbf{turn off} the lights.} \\
\eng{pick up} & recoger & \textit{Can you \textbf{pick up} the phone?} \\
\eng{put on} & ponerse (ropa) & \textit{\textbf{Put on} your jacket.} \\
\eng{take off} & quitarse (ropa) & \textit{\textbf{Take off} your shoes.} \\
\eng{get up} & levantarse & \textit{I \textbf{get up} at 7.} \\
\eng{come back} & volver & \textit{When will you \textbf{come back}?} \\
\eng{find out} & descubrir & \textit{I need to \textbf{find out} the truth.} \\
\eng{wake up} & despertarse & \textit{I \textbf{woke up} late.} \\
\eng{go out} & salir & \textit{Let's \textbf{go out} tonight.} \\
\eng{sit down} & sentarse & \textit{Please \textbf{sit down}.} \\
\eng{stand up} & levantarse & \textit{Please \textbf{stand up}.} \\
\bottomrule
\end{tabularx}
\end{cajachunk}

% ═══════════════════════════════════════════════════════════════════════════════
\section{Oraciones Listas para Minar}
% ═══════════════════════════════════════════════════════════════════════════════

\begin{cajamineria}[Oraciones \textit{i+1} — Semana 7]

\begin{enumerate}[leftmargin=1.2cm]
  \item \textit{I think it \underline{will rain} tomorrow. Take an umbrella.}\\
    {\small\textbf{Objetivo:} futuro con ``will'' para predicciones.}

  \item \textit{We\underline{'re going to visit} my grandmother this weekend.}\\
    {\small\textbf{Objetivo:} futuro con ``going to'' para planes decididos.}

  \item \textit{You \underline{should} drink more water. It's good for you.}\\
    {\small\textbf{Objetivo:} modal ``should'' para consejos.}

  \item \textit{She \underline{can't} come to the party because she has to work.}\\
    {\small\textbf{Objetivo:} modal ``can't'' para imposibilidad.}

  \item \textit{I'm sorry, I \underline{made a mistake} in the exercise.}\\
    {\small\textbf{Objetivo:} la collocation ``make a mistake''.}

  \item \textit{Can you \underline{turn off} the lights, please?}\\
    {\small\textbf{Objetivo:} el phrasal verb ``turn off'' (apagar).}

  \item \textit{He \underline{gave up} smoking last year.}\\
    {\small\textbf{Objetivo:} el phrasal verb ``give up'' (dejar/rendirse).}

  \item \textit{We need to \underline{find out} what happened.}\\
    {\small\textbf{Objetivo:} el phrasal verb ``find out'' (descubrir).}

  \item \textit{You \underline{must} wear a helmet when you ride a bicycle.}\\
    {\small\textbf{Objetivo:} modal ``must'' para obligación.}

  \item \textit{I'll \underline{take a break} and come back in ten minutes.}\\
    {\small\textbf{Objetivo:} la collocation ``take a break'' (tomar un descanso).}
\end{enumerate}
\end{cajamineria}

\vfill
\begin{center}
\rule{0.5\textwidth}{0.5pt}\\[0.3cm]
{\small\color{grissuave} Curso de Inmersión en Inglés · Mes 2 · Semana 7 · Nivel A2}
\end{center}

\end{document}
